\section{Limitations}
\label{sec:limitations}

\paragraph{Scale.}
Our ladder stops at \LargestModel\ parameters. That is a hardware limit, not a
scientific choice: every model runs in float32 on a CPU with 17\,GB of RAM, and
a checkpoint that does not finish enough items in the time available is
excluded rather than reported thin (Appendix~\ref{app:repro}). The trend we
observe is monotone in size, so the honest summary is that we have characterised
the small-model regime and shown the direction of travel, not that we have
located a threshold. Whether frontier-scale models invert at all is untested
here.

\paragraph{Environments.}
ToolShed is our own construction. We built it because the central measurement
needs counterfactual observations that no public benchmark provides, and we have
tried to be explicit about the ways that choice could flatter us: the
perturbation operators are ours, the error messages are ours, and the task
templates are ours. CodeRepair is the check on that, with MBPP problems, real mutations, and
observations written by the Python interpreter rather than by us, and it shows
the same pattern. Neither is a long-horizon, realistic agent
benchmark. We would expect the effect to be larger, not smaller, in settings
with longer transcripts and more failures per trajectory, but we have not shown
that.

\paragraph{Perturbations are synthetic.}
Failing actions come from fixed operators rather than from model samples. This buys the property the scaling claim depends on, that every model is
probed on byte-identical items, at the cost of some realism. Sampling failures from each
model would make the item set model-dependent and would confound size with item
difficulty; we judged that the worse trade. The operators were chosen to match
the mistakes small models actually make, and the rollout study, where models
generate their own failures, is where that choice is put to the test.

\paragraph{Greedy decoding.}
Rollouts use greedy decoding, so each (model, harness) pair yields one
trajectory per task and our intervals over rollouts come from resampling tasks
rather than seeds. Sampling would add a variance component we have not measured.
Greedy is the right default for a comparison of harnesses, since it removes
decoding noise from a contrast we care about, but it is not how every deployed
agent runs.

\paragraph{Log-probabilities of whole strings.}
Corrective gain is a difference in summed token log-probability. Summed
quantities grow with string length, so the raw nat values are not comparable
across environments with very different action lengths. Within an environment
they are comparable across conditions and models because the string being scored
is identical; the normalised repeat probability and the exact greedy repeat rate
are provided precisely because they carry no such caveat.

\paragraph{A numerical failure worth knowing about.}
On this machine, calling \texttt{torch.set\_num\_threads($n$)} with $n > 1$
causes the installed PyTorch build to return all-NaN logits, silently, with
unchanged wall-clock time. We lost one complete probe run to this before
noticing, and our first verification script failed to catch it because
\texttt{max(0.0, nan)} returns \texttt{0.0} in Python, so a comparison of two
NaN quantities reported perfect agreement. Model loading now ends with a forward
pass that must produce finite logits, the scorer checks every returned logit
row, and the verification script rejects non-finite and positive
log-probabilities before comparing them. We record this because it is the kind
of fault that produces a clean-looking table of fabricated numbers, and because the mitigation, which is to assert on the numerical path and not
just on the code path, is cheap and general.

\section{Broader impact}

The interventions we study make small agents repeat themselves less, which
reduces wasted tool calls, wasted tokens and the wall-clock cost of a stuck
loop. The direct effects are mundane and mostly good. Two second-order points
are worth naming. First, a harness that suppresses repeated failed actions makes
an agent more persistent, and persistence is not always desirable: an agent that
keeps trying new calls after several failures may do more damage in a
side-effecting environment than one that gets stuck. Step budgets and
side-effect confirmation remain necessary; nothing here replaces them. Second,
our results make small models more usable as agents, which lowers the cost of
running autonomous systems locally. That is broadly positive for privacy and
access, and it also lowers the cost of running such systems without oversight.
